\begin{papercard}{huckvale2019}
\cf{Stages} \stagetag{A} \stagetag{D} \stagetag{H} (also \stagetag{I}, \stagetag{R}). Policy-vs-traffic study of 36 top-ranked depression and smoking-cessation apps. It is notable for a \emph{personalisation-minimising search protocol}, a pre-specified top-$N$ rule, and \emph{pre-reconciliation} agreement statistics for policy coding.

\cf{Aim} A contemporary assessment of mental-health app privacy practices, specifically whether policy disclosures about third-party transmission match observed behaviour \pc{huckvale2019}{2}{In this study, we aimed to provide a contemporary assessment of the privacy practices}.

\cf{Search protocol} Two authors searched the official Android and iOS stores in the US (14 Jan 2018) and Australia (15 Jan 2018) for ``depression'' and ``smoking cessation''. \emph{Search rank} was used as the popularity proxy, following prior app research \pc{huckvale2019}{2}{We used search rank as a proxy for popularity}. To minimise user-specific tailoring of results, searches ran from an anonymised account with no prior credentials \pc{huckvale2019}{2}{To minimize the risk of user-specific tailoring of search results}. Pre-specified stopping rule: the first 10 apps per term per country-specific store; after pooling and de-duplication, 36 apps (15 Android-only, 14 iOS-only, 7 both) \pc{huckvale2019}{3}{We prespecified that the first 10 apps returned for each search term}. The sampled fraction was estimated against exhaustive prior counts (about 8\,\% of depression apps, 6\,\% of smoking apps). No filter by payment or language \pc{huckvale2019}{3}{Apps were not filtered by payment model or language}. Rank stability was checked at the end of analysis: 72\,\% were still in the top 10 and 92\,\% in the top 20 \pc{huckvale2019}{7}{almost three-quarters (72}.

\cf{Policy capture} A ``permissive stance'' reviewed all policy material presented in or linked from the app UI, linked from the store entry, or on the developer website, including terms and conditions \pc{huckvale2019}{3}{We adopted a permissive stance that reviewed all policy material}.

\cf{Coding scheme} 25 criteria derived from Huckvale 2015, covering primary and secondary uses, transmission, access rights, security, governance, and operational controls such as update procedures. A new item covers procedures after takeover or dissolution of the responsible entity (Table~1) \pc{huckvale2019}{3}{Counts and Proportions of Apps Addressing Specific Privacy Criteria in a Policy}. A distinction was kept between \emph{purposes} of use and \emph{handling by different entities}. Third-party text was coded as advertising/marketing (tailoring adverts or commercial insights about users) vs.\ analytics (tracking use for product improvement) \pc{huckvale2019}{3}{In coding policies, we recognized a distinction between disclosures about the purposes}.

\cf{Reliability} Two independent coders per policy, drawn from three people. Disagreements were resolved by a third person not involved in the initial review. \textbf{Before reconciliation}: raw binary agreement 90.4\,\%, Cohen's $\kappa=0.801$ \pc{huckvale2019}{4}{Raw binary agreement between coders before reconciliation}. The policy reviewer was compensated \pc{huckvale2019}{8}{reviewed the privacy policies}.

\cf{Behavioural test} Two devices (Android 7.1.2, iOS 11.0.1). Two sessions of simulated use exercised the available features, with man-in-the-middle interception including TLS. Destinations and content were tagged automatically and verified manually. A post-hoc review identified social-login functions \pc{huckvale2019}{4}{installed on 1 of 2 test devices}. Reported according to STROBE \pc{huckvale2019}{4}{Strengthening the Reporting of Observational Studies in Epidemiology}.

\cf{Key findings} 69\,\% had a policy. 92\,\% of those mentioned third-party transmission, and transmission was observed in 92\,\% of all apps. 17/36 transmitted to third parties while having no policy (9), not disclosing it (5) or \emph{denying} it (3) \pc{huckvale2019}{4}{Almost half of the apps (17 of 36}. 81\,\% sent data to Google or Facebook, but only 43\,\% and 50\,\% respectively named them \pc{huckvale2019}{5}{Only half of the apps (6 of 12}. Social-login SDKs transmitted identifiers whether or not the login was used \pc{huckvale2019}{5}{Transmissions occurred regardless of}. Linkable identifiers plus the app name could reveal mental-health status, a \emph{capability} for a harm pathway \pc{huckvale2019}{6}{The transmission of even basic details}.

\cf{Limitations stated} A snapshot, because repeated analyses are infeasible given double coding and testing time \pc{huckvale2019}{7}{conducted using a snapshot of apps and policy documentation captured at a single point}. Top 10 only in two areas. Search rank as a heuristic for what users install \pc{huckvale2019}{7}{search rank appears to}. The advertising/analytics categories may be outdated.

\cf{Use for this thesis} Near-template for stage A: anonymised/fresh account, a dated search per storefront, a pre-specified top-$N$ cut per term, pooling and de-duplication, and an end-of-study rank-stability check. For stage I it is a good reporting example: raw agreement \emph{and} $\kappa$ \emph{before} reconciliation, plus a third-person adjudicator. Its ``policy states transmission would not occur'' category is a model for coding \emph{contradicting reassurance} (stance) without measuring traffic. Traffic itself is out of scope.
\end{papercard}
